Alignment intervention studies often report that an intervention ``did not change'' alignment scores.
Such a null is informative only if the measurement could have detected a change, yet alignment evaluations rarely include a positive control.
We build a panel of \emph{synthetic misalignment anchors}, model conditions with a known and independently verified behavioural shift, and measure how often an unmodified battery of eight standard alignment metrics fails to register them.
We test the hypothesis that the false negative rate (\fnr) exceeds 0.30 on two 7B open-weight models with 45 candidate anchors, of which 19 pass verification.
At the battery level the hypothesis is not supported: the pre-specified rule misses 4 of 19 verified anchors (\fnr 0.21, 95\% CI 0.09--0.43), although the interval does not exclude 0.30.
At the single-metric level it is supported: 98 of 152 anchor--metric pairs show no significant shift (\fnr 0.64).
All four battery misses are conditional anchors, whose behaviour is gated on a trigger or hidden from surveys (4 of 7 missed, against 0 of 12 unconditional anchors).
The battery's sensitivity is also not discriminating: the same rule flags 5 of 12 benign system prompts such as ``Use metric units'' as misaligned (false positive rate 0.42).
Sample size matters: at 25 items per metric the battery misses 52\% to 64\% of anchors.
A null from one metric is therefore weak evidence, and a significant shift is not automatically meaningful.
We recommend that intervention studies report a positive control and a benign-perturbation control next to any null.
